% !TeX root = ../../main.tex

Da die Datenbank für jede Transaktion einen SQLSTATE-Code zurückgibt, müssen diese Rückgabewerte entsprechend
ausgewertet werden.
Die Verarbeitung der SQLSTATE-Codes erfolgt im Modul \texttt{retry} innerhalb der Funktion
\texttt{\_start\_transaktion}.
Dort wird die Transaktion in einem \texttt{try: ... except: ...}-Block gestartet.
Tritt während der Ausführung ein von \texttt{psycopg2} Datenbankfehler auf, wird dieser abgefangen und
entsprechend verarbeitet~\ref{lst:errorDetection}.

\begin{lstlisting}[style=pythoncode, caption={Fehlererkennung}, label={lst:errorDetection}]
def _start_transaction():
    try:
        ...
    except psycopg2.Error as e:
        res = "failed"
        pgcode = e.pgcode
def run_retry_with_jitter_delay():
    ...
    while data == "failed" and retries < retry_count:
        ...
\end{lstlisting}

Im Fehlerfall wird neben dem von der Datenbank zurückgegebenen \texttt{pgcode} auch die Variable \texttt{res} auf
den Wert \texttt{failed} gesetzt.
Dadurch wird die weitere Verarbeitung unabhängig vom konkreten SQLSTATE-Code gesteuert.
Der \texttt{pgcode} wird dabei für die spätere Auswertung erfasst, während die Variable \texttt{res} zur Bestimmung des
weiteren Ablaufs des Retry-Mechanismus verwendet wird.

Die eigentliche Entscheidung über einen erneuten Ausführungsversuch erfolgt anschließend in der Retry Funktion.
Dazu wird zunächst geprüft, ob die vorherige Ausführung fehlgeschlagen ist.
Gleichzeitig wird überprüft, ob die maximal zulässige Anzahl an Retries bereits erreicht wurde.
Solange beide Bedingungen erfüllt sind, wird die Transaktion erneut ausgeführt.
Nach jedem fehlgeschlagenen, außer dem initial Versuch wird der Retry-Zähler erhöht, sodass die Anzahl der
erneuten Ausführungen begrenzt bleibt.
Sobald die Transaktion erfolgreich ausgeführt wurde oder die maximale Anzahl an Retries erreicht ist, wird die
Schleife verlassen und das vorliegende Ergebnis zurückgegeben.
Die Begrenzung der Anzahl der Retries ist insbesondere deshalb relevant, da ein Fehler andernfalls zu einer
unbegrenzten Folge von Wiederholungsversuchen führen könnte.

Für die vorliegende Untersuchung ist außerdem zu beachten, dass auf eine dynamische Klassifikation der Fehler
anhand der SQLSTATE-Codes verzichtet wurde.
Die Fehler werden in den Experimenten gezielt erzeugt und dienen jeweils der Untersuchung einer bestimmten
Fehlerklasse.
Stattdessen wird der zurückgegebene SQLSTATE-Code lediglich erfasst.
Im Rahmen der Versuchsauswertung wurde zusätzlich überprüft, dass bei der Durchführung der Experimente
ausschließlich die jeweils vorgesehenen Fehlerklassen auftraten.
Dadurch kann sichergestellt werden, dass die gemessenen Ergebnisse nicht durch unerwartete Fehlerarten beeinflusst
werden.